In recent decades, the seismic assessment of existing buildings has significantly developed from traditional objectives that have focused on ensuring life safety of the building. The economic impact of the 1994 Northridge earthquake in the US due to the extensive damage suffered by buildings, in addition to the overall disruption, highlighted the need for a paradigm shift in the way in which the performance of buildings ought to be defined. The performance of buildings during recent earthquakes in Italy have also emphasised this point where to this day, the city of L’Aquila remains in the process of recovery following the damage induced during the earthquake of 2009, with initial damage from the Rieti earthquake in 2016 showing a similar, if not greater, amount of damage. Such repercussions have highlighted the need for a more comprehensive deliberation of building performance during seismic assessment. Such a methodology that not only considers the life safety of the building occupants during extreme events but also appraises the potential economic consequences of the damage inflicted on the building’s structural and non-structural elements during more frequent events has been developed into what is known as the PEER PBEE methodology. Through this methodology, the performance of buildings can be defined in terms of more meaningful metrics, such as direct economic losses, that can be better understood by a wider audience instead of quantifying the performance through a parlance of demand-capacity checks and verifications mostly familiar to seismic engineers.

This thesis focuses on the assessment of existing Italian reinforced concrete (RC) frame buildings with masonry infill, which were typically gravity load designed (GLD) prior to the introduction of seismic design provisions in the 1970s. The assessment of such GLD RC frames with masonry infill is discussed within a setting similar to that of the FEMA P58 guidelines which aim to provide practising engineers with the tools and procedures, both advanced and simplified, to quantify the performance of existing buildings in a more meaningful way that can be easily conveyed to decision makers. This is applied for these GLD building typologies through the calibration of a numerical modelling procedure in OpenSees to not only account for the non-linearity of the vulnerable components of the structure but also to account for the strength and stiffness degradation so as to appropriately represent the non-linear behaviour and collapse capacity of the building. Additional tools to perform building-specific loss assessment are outlined, where further developments with regard to the proper handling of torsional response, amongst other issues, in the damage and subsequent loss estimation steps of the PEER PBEE are discussed and implemented into a MATLAB based tool. 

Using the numerical modelling approach calibrated to the available experimental test data collected from the literature, the uncertainty associated with the various modelling parameters is quantified via an extensive numerical study. This is done to quantify the propagation of this uncertainty in the modelling of the various structural elements to the overall dispersion of the demand parameters used in damage assessment, in addition to its impact on the collapse fragility function of such buildings. Using these results, empirical values of modelling uncertainty to be used in the assessment of GLD RC frames with masonry infill are quantified. Comparing these empirical values to those available in the literature, it is seen how the fundamental behaviour of the GLD frames differs to the extent that default values available in guidelines such as FEMA P58 cannot be reasonably adopted. Building on this study, extensive numerical analysis is carried out on a variety of case study buildings to quantify and benchmark the performance of such buildings, both in terms of expected demand and overall collapse capacity, where the impact of incorporating the potential shear failure in column members can result in a reduction of up to 20\% of the median collapse intensity due to the interaction with the masonry infill, a failure mechanism that has been highlighted during past earthquakes in Italy. Furthermore, loss estimation studies are carried out on these case study buildings to not only quantify the expected losses but also investigate ways in which shrewd retrofitting can have maximum impact on the overall performance of the buildings. It is shown that when collapse safety of the existing structure is not an issue, as may be the case in areas of low to moderate seismicity in Italy, the performance defined in terms of expected annual loss (EAL) can be effectively reduced through astute retrofitting of the most vulnerable non-structural elements. On the contrary, adopting structural retrofitting schemes involving strengthening and/or stiffening of the structure in compliance with NTC 2008 requirements in these situations may actually lead to a worsening of the building performance defined in terms of EAL. This results due to the trade-off between reduced storey drifts and increased floor accelerations of the retrofitted frame in addition to potential increase in intensity for a given return period due to the shortening of the fundamental period. This observation was also observed in the 2016 earthquake in Ecuador, where excessive strengthening measures to existing hospital structures following the 1998 event resulted in extensive non-structural damage to hospital structures resulting in widespread loss of functionality of these essential services following the 2016 earthquake.

In light of the findings made in relation to the modelling and uncertainty in GLD frames, a new framework for the assessment of GLD RC frames with masonry infill similar to that outlined by FEMA P58 is established. This is described here as a multifaceted framework whereby the extensive approach described above can be adopted, or a more practice-oriented approach that allows simplified structural analysis in the form of displacement-based assessment (DBA) to be adopted. The results using the DBA methodology are compared to the results from refined analysis conducted here to show that although the general framework works quite well, some key ingredients to the DBA methodology are required before a full proposal to GLD RC frames with infills can be made. Nevertheless, the framework is established whereby the median demand on the building can be computed using the completed DBA methodology in conjunction with empirical values of demand dispersion and correlation proposed as part of this study. Comparing this approach with an application to a real school building illustrated both facets of this framework outlined above, where the both the rigorous and simplified approaches are applied to illustrate how the empirical values for demand dispersion and correlations proposed here are to be adopted for such buildings. It is shown how those outlined in FEMA P58 are not representative and tend to underestimate the dispersion, overestimate the demand correlations and subsequently underestimate the expected loss and its associated variance to result in an overall underestimation of the EAL. 

Lastly, the extension of the SAC/FEMA methodology for the computation of the mean annual frequency of exceedance of different limit states is investigated for the structures examined here, and it is shown that the exponent term b required in the demand-intensity relationship should not be taken as unity for GLD RC frames with infill, and instead a set of more appropriate values are proposed for use when masonry infills are present. 
